\documentclass[../main.tex]{subfiles}
\begin{document}

Үйлдвэрлэлийн дадлагыг CallPro Labs байгууллагад 2026 оны 06 сарын
15-наас 07 сарын 02-ны хооронд хийв. Энэ хугацаанд ярианы өгөгдөл
бэлтгэх багийн ажилд оролцож, аудио бичлэгийг текст болгон буулгасан.
Өдөр тутмын ажлын зэрэгцээ дадлагын удирдагчийн зөвлөмжөөр баримтад
тулгуурлан асуултад хариулах RAG чатботын бие даасан төсөл хэрэгжүүлэв.
Тайланд уг төслийн судалгаа, зохиомж, хэрэгжүүлэлт болон туршилтыг
голлон авч үзсэн.

Байгууллагын журам, гарын авлага, тайлан зэрэг баримтаас тодорхой
мэдээлэл олохын тулд хэрэглэгч олон файл нээж, агуулгыг нь унших
шаардлагатай болдог. Түлхүүр үгийн хайлт нь асуулт болон баримтад
өөр үг хэрэглэсэн үед хэрэгтэй хэсгийг олохгүй байж болно. Харин
том хэлний загварт дотоод баримтыг өгөөгүй бол тухайн баримтын
агуулгад тулгуурласан хариулт авах боломжгүй. Иймээс холбогдох
мэдээллийг эхлээд хайж олоод, зөвхөн тэр мэдээллээр хариулт үүсгэх
RAG аргыг сонгосон~\cite{lewis2020rag}.

Төслийн зорилго нь хэрэглэгч өөрийн баримтыг оруулж, асуултын хариултыг
эх сурвалжийн хамт авах боломжтой чатботын ажлын хувилбар бүтээх байв.
Энэ зорилгод хүрэхийн тулд баримтаас текст ялгах, жижиг хэсгүүдэд хуваах,
embedding вектор үүсгэн хадгалах, асуултад тохирох хэсгүүдийг хайх,
олдсон мэдээлэлд тулгуурлан хариулт үүсгэх дарааллыг хэрэгжүүлсэн.
Үүнтэй хамт хэрэглэгчийн бүртгэл, нэвтрэлт, баримтын хандалтын
тусгаарлалт, харилцан ярианы түүхийг хадгалах боломжийг бүрдүүлэв.

Хөгжүүлэлтийг вэб клиент, Fastify API, PostgreSQL болон pgvector
өгөгдлийн сан бүхий бүтэцтэйгээр эхлүүлж, дараа нь мобайл клиент нэмсэн.
Хариулт болон embedding үүсгэхэд Gemini API ашиглав. Интерфэйсийг
монгол хэлээр хийж, кирилл баримт болон латин үсгээр бичсэн асуулттай
ажиллах үед гарсан хайлтын асуудлыг туршилтаар шалгасан.

Ажлын хүрээг PDF-ээс текст ялгах болон текст хуулж оруулах хоёр хэлбэрээр
хязгаарласан. Скан зураг таних, байгууллагын хамтын ажлын талбар,
хариултыг хэсэгчлэн дамжуулах боломжийг энэ хувилбарт хэрэгжүүлээгүй.
Системийг локал орчинд туршсанаас гадна Azure виртуал серверт Docker
Compose-оор байршуулж, Cloudflare Tunnel-ээр HTTPS хаягт холбосон.
Төслийн нээлттэй эх код, баримтжуулалт \url{https://github.com/Tengis01/rag_chatbot}
хаягт, ажиллаж буй вэб системийн туршилтын хувилбар \url{https://ragchatbot.dev}
хаягт нийтэд нээлттэй байршиж байна. GitHub Actions нь кодын шалгалт,
image build-ийг хийж, GHCR дахь хувилбарт image-ээс хяналттай release хийх
дамжлагыг бүрдүүлсэн. Харин хариултын чанарыг өргөн хэмжээний бодит
өгөгдлөөр бүрэн үнэлээгүй. Туршилтын хэсэгт баталгаажуулсан үр дүн болон
шалгаж дуусаагүй ажлыг ялгаж тайлбарласан.

Тайлангийн эхний бүлгүүдэд дадлагын төлөвлөгөө, байгууллага болон
гүйцэтгэсэн ажлын чиг үүргийг танилцуулна. Дараа нь төслийн шаардлага,
ашигласан аргууд, системийн зохиомж, хэрэгжүүлэлтийг тайлбарлаж,
туршилтын үр дүн, гарсан алдаа ба тэдгээрийн шийдлийг өгүүлнэ.
Төгсгөлд эзэмшсэн мэдлэг, ур чадвараа нэгтгэж, байгууллагад өгөх санал
болон цаашдын хөгжүүлэлтийг дүгнэсэн. Хавсралтад өгөгдлийн сангийн
бүтэц, гол функцүүдийн эх код, локал туршилтыг давтах дарааллыг оруулав.

\clearpage
\section*{Нэр томьёоны тодорхойлолт}

Энэхүү тайланд олон улсын програм хангамжийн инженерчлэл, том хэлний
загвар (LLM) болон RAG системийн мэргэжлийн нэр томьёог ашигласан.
Олон утгаар ойлгогдох эрсдэлтэй болон эх хэлээр нь шууд хэрэглэсэн
гол ойлголтуудыг энэ тайланд дараах тодорхой утгаар авч үзэв:

\begin{enumerate}
\setlength{\itemsep}{1.5pt plus 1pt minus 1pt}
\setlength{\parsep}{0pt}
\setlength{\parskip}{0pt}
\item \textbf{Баримтын хэсэг (Chunk).} Оруулсан PDF эсвэл текстийг утга зүйн
дагуу вектор болгон хөрвүүлэхэд тохиромжтой байхаар тодорхой урттай
(500--800 токен) тасдан хуваасан текстийн бие даасан нэгж;
\item \textbf{Вектор дүрслэл (Embedding).} Текстийн утга санааг машин
тооцоолж харьцуулах боломжтой болгон 768 хэмжээст бодит тоон олонлог болгож
буулгасан тоон дүрслэл;
\item \textbf{Эх сурвалжид тулгуурласан хариулт (Grounding).} Хэлний загвар
өөрийн сургагдсан санах ойноос таамаглан зохиож хариулахаас (hallucination)
сэргийлж, зөвхөн олгосон баримтын эх сурвалжийн хүрээнд хариулах инженерийн
зарчим;
\item \textbf{Контекст (Context).} Хэрэглэгчийн асуулттай хамт загварт
дамжуулж буй, өгөгдлийн сангаас хайлтаар олдсон баримтын сонгогдсон
хэсгүүдийн багц;
\item \textbf{Хайлтын босго (Similarity threshold).} Векторын косинус
төстэй байдлын хамгийн бага зөвшөөрөгдөх утга бөгөөд уг босгоос доогуур
оноотой хэсгүүдийг баримтад хамааралгүй гэж үзэж хасна;
\item \textbf{Олон талт хариулт (Maximal Marginal Relevance -- MMR).}
Зөвхөн асуулттай төстэй хэсгүүдийг авахаас гадна өмнө сонгогдсон хэсгүүдтэйгээ
агуулга давхардахгүй, ялгаатай мэдээлэл бүхий баримтын хэсгүүдийг сонгох
эрэмбэлэлтийн арга;
\item \textbf{Хэрэглэгчийн тусгаарлалт (Cross-account isolation).} Нэг
хэрэглэгчийн оруулсан баримт, вектор болон харилцан ярианы түүх нь өөр
хэрэглэгчийн хайлтад хэзээ ч илрэхгүй, систем түвшинд чанд тусгаарлагдах
өгөгдлийн аюулгүй байдал;
\item \textbf{Бүтцийн хувилбаржуулалт (Migration).} Үйлдвэрлэлийн өгөгдлийн
сангийн хүснэгт, баганыг өгөгдөл алдагдуулахгүйгээр дараалсан SQL скриптээр
шат дараатай шинэчлэх үйл явц;
\item \textbf{RAG (Retrieval-Augmented Generation).} Баримтаас мэдээлэл
хайн олж, түүндээ тулгуурлан хариулт үүсгэх том хэлний загварын системчилсэн
архитектур;
\item \textbf{Token (Токен).} Хэлний загварын текстийг боловсруулах
хамгийн бага нэгж. Кирилл монгол үсэг UTF-8 кодчилолд 2 байт эзэлдэг тул
энгийн тэмдэгтийн тооноос ялгаатай хэмжигдэхүүн болдог;
\item \textbf{Drain (Үйлчилгээг түр зогсоох горим).} Серверийг шинэчлэх үед
шинээр ирэх хүсэлтийг түр хязгаарлаж, сервер дээр аль хэдийн эхэлсэн идэвхтэй
боловсруулалтуудыг алдаагүй дуустал хүлээх горим;
\item \textbf{Session.} Хэрэглэгч системд нэвтэрсний дараа төхөөрөмж болон
серверийн хооронд харилцааг баталгаажуулах хугацаат криптограф төлөв;
\item \textbf{HNSW (Hierarchical Navigable Small World).} Олон мянган
вектор дундаас асуултад хамгийн ойр векторуудыг өндөр хурдтай ойролцоолон
олох олон үет граф бүтцийн индекс.
\end{enumerate}

\end{document}
